% Blueprint content for the graphon theory formalization

% Load generated node definitions from LeanArchitect
\input{../../.lake/build/blueprint/library/Graphon}

\chapter{Graphons}

This project formalizes the theory of graphons --- limits of dense graph sequences ---
in Lean~4 using Mathlib, following Lov\'asz~\cite{lovasz2012}.

A \emph{graphon} is a symmetric measurable function $W : [0,1]^2 \to [0,1]$
(more generally, on an arbitrary probability space).  Graphons arise as the
natural limit objects for convergent sequences of dense graphs.

\section{Proof Status}

\textbf{The advertised program is fully proved} (2026-07-10): no live \texttt{sorry}
remains, and every headline theorem verifies with the standard axioms only
(\texttt{propext}, \texttt{Classical.choice}, \texttt{Quot.sound}). The three items
formerly tracked here were each resolved:
\begin{enumerate}
  \item \textbf{Rokhlin / measure-isomorphism input}: the original
    \texttt{exists\_common\_extension} stub was shown unprovable as stated and replaced
    by four corrected cores, all proved from the atomless standard-Borel
    measure-isomorphism theorem, built from scratch in \texttt{Graphon/MeasureIso.lean}
    (mod-0 isomorphism with $([0,1],\lambda)$ plus a mod-0-to-everywhere upgrade); the
    final overlay core is proved in \texttt{Graphon/Overlay.lean}.
    Used by: cut distance triangle inequality, partition alignment, compactness.
  \item \textbf{Algebraic determination} for $k \ge 2$
    (\texttt{matrix\_quotient\_of\_weightedHomSum\_eq}): proved 2026-07-06 via the
    twin-free bijection and the cross-matrix super transfer
    (\texttt{Graphon/CrossSuper.lean}).
  \item \textbf{Determination theorem}
    (\texttt{cutDistance\_zero\_of\_homDensity\_eq}): proved, via the First Sampling
    Lemma (proved 2026-07-08, \texttt{Graphon/SamplingLemma.lean}) and the
    $K$-independent quantitative inverse counting lemma.
\end{enumerate}
Seven \texttt{sorry} declarations are deliberately retained as documentation of
\emph{refuted} conjectures (each marked FALSE/REFUTED in its docstring; none on any
live path).

\section{Main Results}

The formalization establishes:
\begin{itemize}
  \item \textbf{Cut distance pseudometric.}
    Cut distance is a pseudometric: non-negativity, symmetry
    (Proposition~\ref{thm:cutDistance-symm}), and the triangle inequality
    (Theorem~\ref{thm:cutDistance-triangle}).
  \item \textbf{Frieze--Kannan weak regularity lemma.}
    Every graphon admits a step approximation of bounded complexity
    (Theorem~\ref{thm:regularity}), following
    Frieze--Kannan~\cite{friezekannan1999} and
    Lov\'asz~\cite[Corollary~9.13]{lovasz2012}.
  \item \textbf{Counting lemma.}
    Small cut distance implies similar homomorphism densities
    (Theorem~\ref{thm:counting-lemma}).
  \item \textbf{Compactness.}
    The cut-distance quotient modulo weak isomorphism is a compact metric
    space; concretely, the graphon pseudometric space is totally bounded
    (Theorem~\ref{thm:totallyBounded}) and complete
    (Theorem~\ref{thm:complete}).
  \item \textbf{Inverse counting lemma.}
    For any $\varepsilon > 0$, finitely many test graphs control cut distance
    up to $\varepsilon$ (Theorem~\ref{thm:inverse-counting}),
    following Lov\'asz~\cite[Lemma~10.32]{lovasz2012}.
  \item \textbf{Convergence equivalence.}
    Cut distance convergence is equivalent to convergence of all
    homomorphism densities (Theorem~\ref{thm:convergence-equiv}).
\end{itemize}

\section{Graphon Definition}

A graphon on a probability space $(\alpha, \mu)$ is a symmetric measurable
function taking values in $[0,1]$ a.e., following
Lov\'asz~\cite[Definition~7.1]{lovasz2012}.

\inputleannode{def:graphon}


\chapter{Cut Distance}

The \emph{cut distance} is the fundamental pseudometric on graphon space, measuring
how close two graphons are after optimal rearrangement of the underlying
probability space.  The development follows
Lov\'asz~\cite[Chapter~8]{lovasz2012} and
Borgs--Chayes--Lov\'asz--S\'os--Vesztergombi~\cite{borgs2008}.

\section{Cut Norm}

The cut norm of the difference of two graphons, defined as the supremum
of $|\int_{S \times T} (U - W)|$ over measurable sets $S, T$.

\inputleannode{def:cutNormDiff}

\section{Pullback}

The pullback of a graphon under a measure-preserving map, fundamental to
the definition of cut distance.  Following
Lov\'asz~\cite[Section~8.2]{lovasz2012}.

\inputleannode{def:pullback}

\section{Cut Distance Definition}

The cut distance between two graphons, defined as the infimum of
$\|U^\varphi - W^\psi\|_\square$ over measure-preserving maps $\varphi, \psi$.
This two-sided formulation avoids the need for invertibility
(Lov\'asz~\cite[Section~8.2]{lovasz2012}).

\inputleannode{def:cutDistance}

\section{Metric Properties}

Cut distance is a pseudometric on graphon space
(Lov\'asz~\cite[Theorem~8.13, Corollary~8.14]{lovasz2012}).

\inputleannode{thm:cutDistance-nonneg}

\inputleannode{thm:cutDistance-self}

\inputleannode{thm:cutDistance-symm}

The triangle inequality uses Rokhlin's theorem (that standard Borel
probability spaces are isomorphic) to align measure-preserving witnesses.
\emph{Status:} proved modulo a \texttt{sorry}'d Rokhlin interface
(\texttt{exists\_common\_extension}), which is planned to be proved in future work.

\inputleannode{thm:cutDistance-triangle}

Pullback by a measure-preserving bijection preserves cut distance
(Lov\'asz~\cite[Section~8.2]{lovasz2012}).

\inputleannode{thm:cutDistance-pullback-eq-zero}


\chapter{Partitions and Step Graphons}

Step graphons --- graphons that are constant on rectangles of a partition ---
are the finite-dimensional building blocks of graphon theory.

\section{Measurable Partitions}

A finite partition of the probability space into measurable sets that
cover almost everywhere.

\inputleannode{def:measurablePartition}

\section{Stepification}

The stepification $W_P$ of a graphon $W$ with respect to a partition $P$
replaces $W$ on each rectangle $S \times T$ by its average
(Lov\'asz~\cite[Section~9.2]{lovasz2012}).

\inputleannode{def:stepify}

\section{Step Graphons from Coefficients}

A step graphon built from explicit coefficients on a partition.

\inputleannode{def:mkStepGraphon}


\chapter{Regularity}

The Frieze--Kannan weak regularity lemma for graphons: every graphon can be
approximated by a step graphon of bounded complexity.
Following Frieze--Kannan~\cite{friezekannan1999} and
Lov\'asz~\cite[Corollary~9.13]{lovasz2012}.

\section{Energy Increment}

The energy increment lemma: if a partition $P$ does not approximate $W$ well,
then a refinement $Q$ has strictly higher energy
(Lov\'asz~\cite[Lemma~9.11]{lovasz2012}).

\inputleannode{thm:energy-increment}

\section{Regularity Lemma}

The regularity lemma: for any $\varepsilon > 0$, every graphon admits a
partition $P$ with at most $\mathrm{regularityBound}(\varepsilon)$ parts such that
$\|W - W_P\|_\square \le \varepsilon$.

\inputleannode{thm:regularity}


\chapter{Homomorphism Densities and the Counting Lemma}

Homomorphism densities encode the local structure of a graphon.
The counting lemma shows that small cut distance implies similar
homomorphism densities.  Following
Lov\'asz~\cite[Chapter~5 and Section~10.1]{lovasz2012}.

\section{Homomorphism Density}

For a finite graph $F$ and a graphon $W$, the homomorphism density
$t(F, W) = \int \prod_{uv \in E(F)} W(x_u, x_v) \, d\mu^{V(F)}$
(Lov\'asz~\cite[Equation~(5.29)]{lovasz2012}).

\inputleannode{def:homDensity}

\section{Counting Lemma}

The counting lemma: $|t(F, U) - t(F, W)| \le |E(F)| \cdot \|U - W\|_\square$
(Lov\'asz~\cite[Lemma~10.23]{lovasz2012}).

\inputleannode{thm:counting-lemma}


\chapter{Compactness}

The cut-distance quotient modulo weak isomorphism is a compact metric space
(Lov\'asz~\cite[Theorem~9.23]{lovasz2012}).  Concretely, we prove total
boundedness (via the regularity lemma and grid quantization) and
completeness (via a direct limit construction from rapidly converging
subsequences) of the graphon pseudometric space.

\section{Total Boundedness}

For any $\varepsilon > 0$, there exists a finite $\varepsilon$-net in cut distance.
The construction uses the regularity lemma to approximate any graphon by
a step graphon, then quantizes coefficients to a grid
(Lov\'asz~\cite[Section~9.3]{lovasz2012}).

\inputleannode{thm:totallyBounded}

\section{Completeness}

Every Cauchy sequence in cut distance has a limit graphon.
The proof extracts a rapidly converging subsequence, builds the limit
graphon via Radon--Nikodym, and shows the full sequence converges
(Lov\'asz~\cite[Section~9.3]{lovasz2012}).

\inputleannode{thm:complete}


\chapter{Inverse Counting Lemma and Convergence}

The inverse counting lemma is the converse of the counting lemma:
if all homomorphism densities are similar, then the graphons are close
in cut distance.  Together with the counting lemma, this gives the
fundamental convergence equivalence
(Lov\'asz~\cite[Section~10.6 and Theorem~11.5]{lovasz2012}).

\section{Inverse Counting Lemma}

For any $\varepsilon > 0$, there exist $\delta > 0$ and a finite family of
test graphs such that $\delta$-close homomorphism densities imply
$\varepsilon$-close cut distance
(Lov\'asz~\cite[Lemma~10.32]{lovasz2012}).
\emph{Status:} proved modulo both temporary \texttt{sorry}s (Rokhlin for partition
alignment, algebraic determination for the step graphon core).

\inputleannode{thm:inverse-counting}

\section{Convergence Equivalence}

A sequence of graphons converges in cut distance if and only if all
homomorphism densities converge
(Lov\'asz~\cite[Theorem~11.5]{lovasz2012}).

\inputleannode{thm:convergence-equiv}

\bibliographystyle{amsalpha}
\bibliography{refs}
